\chapter{AI Hat Setup and TPU Verification}

This chapter documents the preparatory steps required to enable and verify the M.2 AI Hat on the Raspberry Pi 5. 
While the AI Hat was physically installed as part of the system hardware, hardware-accelerated TPU inference was 
not utilized in the final deployed pipeline. The following steps therefore represent optional validation and future 
extension capabilities rather than the active inference configuration used in this project.

\section{Verifying TPU Detection}
Detection of the Edge TPU can be verified at the operating system level using standard Linux utilities:
\begin{verbatim}
	lsusb | grep -i google
	dmesg | grep -i edgetpu
\end{verbatim}

These commands confirm whether the AI Hat is recognized by the system kernel and PCIe subsystem. 
In this project, inference was performed using CPU-based TensorFlow Lite execution.

\section{Installing TPU Libraries}
Support libraries for the Edge TPU can be installed using the system package manager:
\begin{verbatim}
	sudo apt install python3-edgetpu python3-pycoral -y
\end{verbatim}

These libraries are required only when deploying models compiled specifically for TPU execution and were not required 
for the final implementation.

\section{Running TPU Inference}
A sample command for validating TPU-based inference is shown below for completeness:
\begin{verbatim}
	python3 -m pycoral.examples.classify_image \
	--model mobilenet_v2_edgetpu.tflite \
	--labels imagenet_labels.txt \
	--input test.jpg
\end{verbatim}

This example demonstrates standalone TPU inference using a precompiled Edge TPU model. 
The deployed bird recognition system instead used a standard TensorFlow Lite model executed on the Raspberry Pi CPU, 
ensuring compatibility with the implemented Python inference pipeline.

